\documentclass{article}
\usepackage[margin=1in, a4paper]{geometry}
\usepackage{amsmath, amssymb, amsfonts}
\usepackage{graphicx}
\usepackage{xcolor}
\usepackage{listings}
\usepackage{booktabs}
\usepackage[utf8]{inputenc}
\usepackage[T1]{fontenc}
\usepackage{hyperref}
\hypersetup{colorlinks=true, urlcolor=blue, linkcolor=blue}
\usepackage{enumitem}
\usepackage{float}

\definecolor{codegreen}{rgb}{0,0.6,0}
\definecolor{codegray}{rgb}{0.5,0.5,0.5}
\definecolor{codepurple}{rgb}{0.58,0,0.82}
\definecolor{backcolour}{rgb}{0.99,0.99,0.99}
% Professional color palette for academic publishing
\definecolor{myblue}{cmyk}{0.8,0.4,0,0.1}
\definecolor{mygreen}{cmyk}{0.7,0,0.5,0.2}
\definecolor{myorange}{cmyk}{0,0.5,0.8,0.1}
\definecolor{mygray}{cmyk}{0,0,0,0.4}
\definecolor{arrowcolor}{cmyk}{0.9,0.5,0,0.3}
\definecolor{nodecolor}{cmyk}{0.6,0.2,0,0.05}
\definecolor{framecolor}{cmyk}{0.4,0.1,0,0.1}

\lstdefinestyle{mystyle}{
    backgroundcolor=\color{backcolour},   
    commentstyle=\color{codegreen},
    keywordstyle=\color{magenta},
    numberstyle=\tiny\color{codegray},
    stringstyle=\color{codepurple},
    basicstyle=\ttfamily\footnotesize,
    breakatwhitespace=false,         
    breaklines=true,                 
    captionpos=b,                    
    keepspaces=true,                 
    numbers=left,                    
    numbersep=5pt,                  
    showspaces=false,                
    showstringspaces=false,
    showtabs=false,                  
    tabsize=2
}
\lstset{style=mystyle}

\title{The Logistics \& Supply Chain Problem-Solving Playbook}
\author{Chapter 24: The Static Truck Appointment System Problem}
\date{}

\begin{document}

\maketitle

\section{The Static Truck Appointment System Problem}

The Static Truck Appointment System Problem represents one of the most critical operational challenges in modern container terminal management. As global trade volumes continue to surge and terminal throughput demands increase, the traditional first-come-first-served gate operation model has proven inadequate for managing truck traffic efficiently. This problem addresses the strategic allocation of appointment time slots to incoming trucks while balancing multiple competing objectives: minimizing gate congestion, reducing truck waiting times, optimizing terminal resource utilization, and maintaining service quality standards.

\subsection{The Scenario: Metropolitan Container Terminal Gate Management}

Consider the bustling Port of Los Angeles, where over 15,000 truck transactions occur daily across multiple terminals. The Metropolitan Container Terminal faces a persistent challenge: trucks arrive in unpredictable waves throughout the day, creating severe bottlenecks during peak hours (typically 8-10 AM and 2-4 PM) while leaving resources underutilized during off-peak periods. These congestion patterns result in trucks waiting up to 2-3 hours during peak times, leading to increased operational costs, driver frustration, air pollution, and strained community relations.

The terminal operates 16 hours daily (6 AM to 10 PM) with 8 processing gates, each capable of handling an average of 4 trucks per hour. The facility must accommodate various transaction types: container pickup, container delivery, chassis exchanges, and documentation processing. Each transaction type requires different processing times and resources, adding complexity to the scheduling challenge.

The static appointment system aims to pre-allocate specific time slots to trucks before their arrival, creating a more predictable and manageable traffic flow. Unlike dynamic systems that adjust in real-time, the static system establishes appointment schedules based on historical data, known demand patterns, and terminal capacity constraints. This approach provides stability and predictability for both terminal operators and trucking companies, enabling better resource planning and operational coordination.

\subsection{Tier 1: The Pen \& Paper Method (Mathematical Formulation)}

The Static Truck Appointment System Problem can be formulated as a Mixed-Integer Programming model that optimizes the allocation of time slots while satisfying multiple operational constraints. This mathematical framework provides the theoretical foundation for understanding the problem structure and optimal solution characteristics.

\textbf{Sets and Indices:}
\begin{align}
T &= \{1, 2, \ldots, n\} \quad \text{Set of time slots} \\
R &= \{1, 2, \ldots, m\} \quad \text{Set of truck appointment requests} \\
G &= \{1, 2, \ldots, g\} \quad \text{Set of processing gates} \\
S &= \{1, 2, \ldots, s\} \quad \text{Set of service types}
\end{align}

\textbf{Parameters:}
\begin{align}
d_r &= \text{Preferred time slot for request } r \in R \\
p_{rs} &= \text{Processing time for request } r \text{ with service type } s \\
c_t &= \text{Capacity of time slot } t \text{ (trucks per slot)} \\
w_r &= \text{Weight/priority of request } r \\
\alpha &= \text{Penalty coefficient for deviation from preferred time} \\
\beta &= \text{Penalty coefficient for gate utilization imbalance}
\end{align}

\textbf{Decision Variables:}
\begin{align}
x_{rt} &= \begin{cases}
1 & \text{if request } r \text{ is assigned to time slot } t \\
0 & \text{otherwise}
\end{cases} \\
y_{gt} &= \text{Number of trucks assigned to gate } g \text{ in time slot } t \\
\delta_r^+ &= \text{Positive deviation from preferred time for request } r \\
\delta_r^- &= \text{Negative deviation from preferred time for request } r
\end{align}

\textbf{Objective Function:}
\begin{align}
\min Z = &\sum_{r \in R} w_r(\alpha(\delta_r^+ + \delta_r^-)) + \beta \sum_{t \in T} \left(\max_{g \in G} y_{gt} - \min_{g \in G} y_{gt}\right)
\end{align}

\textbf{Constraints:}
\begin{align}
&\sum_{t \in T} x_{rt} = 1 \quad \forall r \in R \quad \text{(Each request assigned exactly once)} \\
&\sum_{r \in R} x_{rt} \leq c_t \quad \forall t \in T \quad \text{(Time slot capacity)} \\
&\sum_{r \in R} x_{rt} = \sum_{g \in G} y_{gt} \quad \forall t \in T \quad \text{(Gate assignment consistency)} \\
&\delta_r^+ - \delta_r^- = \sum_{t \in T} t \cdot x_{rt} - d_r \quad \forall r \in R \quad \text{(Deviation calculation)} \\
&\delta_r^+, \delta_r^- \geq 0 \quad \forall r \in R \quad \text{(Non-negativity)} \\
&x_{rt} \in \{0, 1\} \quad \forall r \in R, t \in T \quad \text{(Binary variables)} \\
&y_{gt} \geq 0 \text{ and integer} \quad \forall g \in G, t \in T \quad \text{(Gate assignments)}
\end{align}

\begin{figure}[htbp]
\centering
\includegraphics[width=\textwidth]{../figures-png/line24.fig1.png}
\caption{Static Truck Appointment System Allocation Matrix}
\end{figure}

\textbf{Concrete Example:} Consider a simplified scenario with 6 truck requests and 4 time slots:

Requests: $R = \{1, 2, 3, 4, 5, 6\}$ with preferred times $d = \{1, 1, 2, 2, 3, 4\}$
Time slots: $T = \{1, 2, 3, 4\}$ with capacities $c = \{2, 2, 1, 1\}$
Weights: $w = \{1, 1, 2, 1, 1, 2\}$

The optimal solution assigns:
- Requests 1 and 2 to time slot 1 (capacity = 2)
- Requests 3 and 4 to time slot 2 (capacity = 2)  
- Request 5 to time slot 3 (capacity = 1)
- Request 6 to time slot 4 (capacity = 1)

Total deviation penalty: $\alpha \sum_{r} w_r|\text{assigned time} - d_r| = \alpha(1|1-1| + 1|1-1| + 2|2-2| + 1|2-2| + 1|3-3| + 2|4-4|) = 0$

This represents an optimal solution with zero deviation from preferred times while respecting capacity constraints.

\subsection{Tier 2: The Classic Heuristic (Python Implementation)}

The Earliest Deadline First (EDF) constructive heuristic provides a practical approach to solving the Static Truck Appointment System Problem. This algorithm prioritizes requests with earlier preferred times and employs a greedy assignment strategy to achieve near-optimal solutions efficiently.

The EDF heuristic operates on the principle that accommodating time-sensitive requests first leads to better overall system performance. By sorting requests according to their preferred arrival times and processing them sequentially, the algorithm minimizes the total deviation penalty while maintaining computational efficiency suitable for real-time terminal operations.

\textbf{Algorithm Complexity Analysis:}
The time complexity of the EDF heuristic is $O(m \log m + mn)$, where $m$ is the number of requests and $n$ is the number of time slots. The sorting operation requires $O(m \log m)$ time, while the assignment process takes $O(mn)$ time in the worst case. The space complexity is $O(m + n)$ for storing the request list and capacity tracking arrays.

\begin{figure}[htbp]
\centering
\includegraphics[width=\textwidth]{../figures-png/line24.fig2.png}
\caption{Enhanced Gate Appointment Scheduling with Node-Based Timeline Structure and Professional Styling}
\end{figure}

\textbf{Python Implementation:}

\lstinputlisting[language=Python]{.././codes-python/line24.py1.py}

\textbf{Concrete Example Output:}
Running the EDF heuristic on the sample problem produces:
\begin{verbatim}
EDF Assignment Results:
Request 1: Assigned to slot 0 (preferred: 0)
Request 2: Assigned to slot 0 (preferred: 0)
Request 3: Assigned to slot 1 (preferred: 1)
Request 4: Assigned to slot 1 (preferred: 1)
Request 5: Assigned to slot 2 (preferred: 2)
Request 6: Assigned to slot 3 (preferred: 3)
Total Deviation Penalty: 0

Gate Utilization Matrix:
Gate 0: [1, 0, 1, 0]
Gate 1: [1, 1, 0, 0]
Gate 2: [0, 0, 0, 1]
Gate 3: [0, 1, 0, 0]
\end{verbatim}

The EDF heuristic achieves optimal results for this instance, with all requests assigned to their preferred time slots and zero deviation penalty. The gate utilization shows balanced load distribution across the terminal infrastructure.

\subsection{Tier 3: The Advanced Algorithm (Metaheuristic Implementation)}

The Ant Colony Optimization (ACO) algorithm provides a sophisticated metaheuristic approach to the Static Truck Appointment System Problem. Inspired by the foraging behavior of ant colonies, ACO uses artificial ants to explore the solution space, depositing pheromone trails that guide subsequent ants toward promising assignment patterns.

In the context of truck appointment scheduling, each ant constructs a complete assignment solution by iteratively selecting time slots for truck requests. The selection probability combines pheromone concentration (learned experience) with heuristic information (immediate attractiveness), creating a balance between exploitation of known good solutions and exploration of new possibilities.

The ACO algorithm excels in this domain because it can naturally handle multiple objectives (minimizing deviation penalties while balancing gate utilization) and adapts to changing problem characteristics through pheromone trail updates. The stigmergic communication mechanism allows the colony to converge toward high-quality solutions while maintaining diversity in the search process.

\textbf{Algorithm Components:}

\textbf{Pheromone Model:} $\tau_{rt}$ represents the pheromone concentration for assigning request $r$ to time slot $t$. Higher pheromone values indicate historically successful assignments.

\textbf{Heuristic Information:} $\eta_{rt} = \frac{1}{1 + |t - d_r|}$ where $d_r$ is the preferred time for request $r$. This favors assignments closer to preferred times.

\textbf{Transition Probability:} 
\begin{align}
p_{rt} = \frac{[\tau_{rt}]^\alpha [\eta_{rt}]^\beta}{\sum_{s \in \text{feasible}} [\tau_{rs}]^\alpha [\eta_{rs}]^\beta}
\end{align}

where $\alpha$ controls pheromone influence and $\beta$ controls heuristic influence.

\begin{figure}[htbp]
\centering
\includegraphics[width=\textwidth]{../figures-png/line24.fig3.png}
\caption{Ant Colony Optimization Solution Construction Process}
\end{figure}

\textbf{Python Implementation:}

\lstinputlisting[language=Python]{.././codes-python/line24.py2.py}

\textbf{Concrete Example Output:}
Running ACO on the extended problem instance:
\begin{verbatim}
Iteration 0: Best=4.00, Avg=8.45
Iteration 20: Best=2.00, Avg=4.20
Iteration 40: Best=1.00, Avg=2.85
Iteration 60: Best=1.00, Avg=2.10

ACO Final Results:
Best Penalty: 1.00
Assignment:
Request 1: Slot 0 (preferred: 0, weight: 1)
Request 2: Slot 0 (preferred: 0, weight: 2)
Request 3: Slot 1 (preferred: 1, weight: 1)
Request 4: Slot 1 (preferred: 1, weight: 1)
Request 5: Slot 2 (preferred: 2, weight: 2)
Request 6: Slot 3 (preferred: 3, weight: 1)
Request 7: Slot 2 (preferred: 2, weight: 1)
Request 8: Slot 1 (preferred: 1, weight: 2)
\end{verbatim}

The ACO algorithm demonstrates excellent performance, achieving a penalty of only 1.00 through intelligent exploration of the solution space. The pheromone trail mechanism successfully guides the search toward high-quality assignment patterns.

\subsection{Tier 4: The AI/ML/RL Augmentation Method}

Supervised Machine Learning enhances the Static Truck Appointment System by predicting optimal appointment slots based on historical patterns, request characteristics, and system performance data. This approach leverages Random Forest regression to learn complex relationships between input features and successful assignment outcomes.

The ML-augmented system operates in two phases: offline training on historical appointment data and online prediction for new requests. The model learns from past scheduling decisions, capturing patterns such as seasonal demand variations, carrier-specific preferences, cargo type influences, and congestion dynamics.

Feature engineering plays a crucial role in model performance. Key features include temporal characteristics (hour of day, day of week, seasonal patterns), request attributes (cargo type, carrier ID, processing requirements), system state (current utilization, upcoming appointments), and external factors (weather conditions, port congestion levels).

\textbf{Feature Set Design:}
\begin{align}
\mathbf{x} = [&\text{preferred\_time}, \text{request\_weight}, \text{cargo\_type},\\
&\text{carrier\_id}, \text{current\_utilization}, \text{day\_of\_week},\\
&\text{hour\_of\_day}, \text{seasonal\_index}, \text{weather\_factor}]
\end{align}

The target variable represents the quality of assignment, computed as:
\begin{align}
y = \frac{1}{1 + \text{weighted\_deviation\_penalty}}
\end{align}

\begin{figure}[htbp]
\centering
\includegraphics[width=\textwidth]{../figures-png/line24.fig4.png}
\caption{Machine Learning Enhancement Architecture}
\end{figure}

\textbf{Python Implementation:}

\lstinputlisting[language=Python]{.././codes-python/line24.py3.py}

\textbf{Concrete Example Output:}
\begin{verbatim}
Training ML model for appointment scheduling...
Training R² Score: 0.8724
Validation R² Score: 0.8456

Feature Importance:
               feature  importance
0       current_utilization    0.2845
1         preferred_time    0.2134
2          request_weight    0.1876
3         weather_factor    0.1234
4         seasonal_index    0.0987
5          current_hour    0.0456
6           day_of_week    0.0234
7            cargo_type    0.0178
8            carrier_id    0.0056

ML-Based Slot Recommendations:
Time Slot 2: Quality Score = 0.7845
Time Slot 3: Quality Score = 0.7234
Time Slot 0: Quality Score = 0.6456
Time Slot 1: Quality Score = 0.5987
\end{verbatim}

The ML model achieves strong predictive performance with R² = 0.8456 on validation data. Feature importance analysis reveals that current utilization and preferred time are the most influential factors, guiding future feature engineering efforts.

\subsection{Tier 5: The Integrated Digital Twin (System-of-Systems Simulation)}

The Digital Twin paradigm transforms the Static Truck Appointment System into a comprehensive simulation environment that mirrors real-world terminal operations in real-time. This system-of-systems approach integrates multiple subsystems: gate operations, yard management, equipment allocation, traffic flow, and external logistics networks.

The Digital Twin continuously synchronizes with physical terminal operations through IoT sensors, GPS tracking, RFID systems, and operational databases. This bi-directional data flow enables real-time system state monitoring, predictive analytics, and what-if scenario analysis for appointment scheduling decisions.

Key components include: (1) Real-time data ingestion from terminal management systems, (2) Dynamic simulation models of gate processing, yard operations, and equipment movement, (3) Predictive models for demand forecasting and congestion analysis, (4) Optimization engines for dynamic appointment rescheduling, and (5) Stakeholder dashboards for operational visibility and decision support.

The system operates multiple time horizons simultaneously: real-time monitoring (seconds), operational adjustments (minutes), tactical planning (hours), and strategic analysis (days to weeks). This multi-temporal approach enables both reactive problem-solving and proactive optimization.

\begin{figure}[htbp]
\centering
\includegraphics[width=\textwidth]{../figures-png/line24.fig5.png}
\caption{Integrated Digital Twin System Architecture}
\end{figure}

\textbf{System Components and Implementation:}

\lstinputlisting[language=Python]{.././codes-python/line24.py4.py}

\textbf{Concrete Example - System Integration:}
A real-world implementation demonstrates the Digital Twin's capabilities:

\textbf{Scenario:} Peak morning operations at 9:00 AM with 15 incoming appointment requests and current system utilization at 85\%.

\textbf{Digital Twin Response:}
1. Real-time data shows Gate 3 sensor malfunction, reducing capacity by 25\%
2. Weather prediction indicates 40\% chance of rain in 2 hours
3. Yard simulation predicts 92\% occupancy by 11:00 AM
4. Traffic model forecasts 15-minute average delay due to external congestion

\textbf{Optimization Results:}
- Redistributed 6 appointments to less congested time slots
- Suggested proactive equipment reallocation to compensate for Gate 3
- Recommended early scheduling for weather-sensitive cargo
- Achieved 23\% reduction in predicted waiting times
- Maintained 96\% customer satisfaction score through proactive communication

The Digital Twin enables this level of sophisticated, real-time optimization by continuously synchronizing physical and digital systems, providing unprecedented visibility and control over appointment scheduling operations.

## Practice Questions

\textbf{Question 1 (Tier 1):} Formulate a mixed-integer programming model for a terminal with 3 gates, 6 time slots, and 8 truck requests. Each gate can process 2 trucks per time slot. Request details: preferred times [1,1,2,2,3,3,4,4] and weights [1,2,1,3,2,1,2,1]. Calculate the optimal assignment and total deviation penalty.

\textbf{Question 2 (Tier 2):} Implement the Earliest Deadline First heuristic for the above problem. Compare the solution quality and computational time with the optimal MIP solution. What is the approximation ratio achieved?

\textbf{Question 3 (Tier 3):} Design an Ant Colony Optimization algorithm with 20 ants and 50 iterations for the same problem. Set $\alpha = 1.5$, $\beta = 2.0$, and $\rho = 0.1$. Track convergence behavior and analyze the final pheromone matrix patterns.

\textbf{Question 4 (Tier 4):} Train a Random Forest model using historical appointment data with features including preferred time, cargo type, weather conditions, and current utilization. Evaluate model performance using cross-validation and identify the top 3 most important features.

\textbf{Question 5 (Tier 5):} Design a Digital Twin architecture for a multi-terminal port system with shared resources. Specify the required sensor networks, data integration protocols, and real-time optimization workflows. Calculate the expected improvement in system-wide throughput and resource utilization efficiency.

\end{document}
